\chapter{Discussion}
\label{chap:discussion}

% A chapter outline, one clause per section, in the style of the thesis outline
% in the introduction. It says how the chapter is built, not what it argues.
\Cref{sec:disc-reading} reads the main result against the baseline, \Cref{sec:disc-mechanism} traces it to a mechanism, and \Cref{sec:limitations} closes with the limitations and the threats to validity.

\section{What the Numbers Support}
\label{sec:disc-reading}

Sed ut perspiciatis unde omnis iste natus error sit voluptatem accusantium doloremque laudantium, totam rem aperiam, eaque ipsa quae ab illo inventore veritatis (\Cref{tab:exp-headline}).
A discussion interprets results that a previous chapter established, so it cites them rather than restating the table.
The gap between 0.501 and 0.909 holds under every fold set, which is what makes it a finding instead of an artifact of one split.

\section{Why It Happens}
\label{sec:disc-mechanism}

Nemo enim ipsam voluptatem quia voluptas sit aspernatur aut odit aut fugit, sed quia consequuntur magni dolores eos qui ratione voluptatem sequi nesciunt.
Neque porro quisquam est, qui dolorem ipsum quia dolor sit amet, consectetur, adipisci velit, sed quia non numquam eius modi tempora incidunt ut labore et dolore magnam aliquam quaerat voluptatem~\cite{example_JournalArticle_2023}.

\paragraph{A named step in the argument.}
Ut enim ad minima veniam, quis nostrum exercitationem ullam corporis suscipit laboriosam, nisi ut aliquid ex ea commodi consequatur.
A paragraph heading carries one step of the argument and keeps a long section navigable.

% The limitations are a section of this chapter, \input from their own file so that
% they can be written and reviewed separately. \input does not force a page break,
% so the section follows the text above on the same page.
\section{Limitations and Threats to Validity}
\label{sec:limitations}

% A section, not a chapter: it is \input from chapters/discussion.tex.
Every result rests on one corpus, one protocol, and three seeds, so each number carries the uncertainty of that setup.
At vero eos et accusamus et iusto odio dignissimos ducimus qui blanditiis praesentium voluptatum deleniti atque corrupti quos dolores et quas molestias excepturi sint.

\paragraph{Construct validity.}
Et harum quidem rerum facilis est et expedita distinctio, and the measure may capture something adjacent to what the question asks about.

\paragraph{Internal validity.}
Nam libero tempore, cum soluta nobis est eligendi optio cumque nihil impedit quo minus id quod maxime placeat facere possimus.

\paragraph{External validity.}
Temporibus autem quibusdam et aut officiis debitis aut rerum necessitatibus saepe eveniet ut et voluptates repudiandae sint et molestiae non recusandae.

